\section{DOS PODERES DE REPRESENTAÇÃO}

\begin{enumerate}
\item O subscritor comprova a legitimidade para representar o Município mediante procuração regularmente outorgada pelo(a) Prefeito(a), juntada no Anexo I, suficiente para fins de atuação administrativa perante a ANEEL e a distribuidora, nos termos do art. 9º da REN nº 1.000/2021, da Lei nº 9.784/1999 (arts. 2º e 4º) e, principalmente, do Estatuto da Advocacia – Lei nº 8.906/1994.
\item A exigência de quaisquer documentos além da procuração (como contrato de mandato/contrato de honorários) carece de amparo legal e viola as prerrogativas profissionais do advogado (Lei nº 8.906/1994). 
\item Sem reconhecimento de obrigação e para evitar dilações indevidas, apresenta-se, por liberalidade, cópia do contrato de prestação de serviços firmado com o Município (reclamante).
\item Esclarece-se, por fim, que o Município celebrou contrato com a sociedade profissional do Dr. Abel Gomes Cunha, atribuindo-lhe poderes específicos para representação perante a ANEEL e a distribuidora, conforme cláusula contratual cuja íntegra segue no Anexo II. Requer-se, portanto, o pleno reconhecimento da regularidade da representação nos autos.

\Citacao{\clausula{}}

\item O referido item contém cláusula expressa de outorga de poderes de representação processual ao representante do escritório contratado pelo Município, nos termos do Contrato de Mandato já acostado, tornando incontestável a legitimidade para atuar em nome do Município perante esta distribuidora.
\item Ressalte-se que o referido contrato foi regularmente publicado no Diário Oficial do Município, conforme exige a legislação vigente, circunstância que assegura sua plena eficácia jurídica e produz os efeitos legais inerentes ao instrumento contratual. Portanto, o requisito formal de publicação do contrato foi regulamente atendido. 
\end{enumerate}
